\chapter{A Native PhaseNav Theta-Bridge Construction}
\label{chap:phasenav-native}

\section{Purpose and claim boundary}

This chapter supplies the first executable construction aimed at the open bridge
identified in \cref{chap:candidates}.  The construction is expressed natively as
a $36$-dimensional weighted phase state rather than as prose subsequently encoded
into PhaseNav.  Its source of truth is the program
\begin{center}
\texttt{construction/phasenav/secret\_of\_half\_theta\_bridge.pnv}.
\end{center}
The accompanying Python module parses the node basis and execution profile from
that file and evaluates the state.  It is therefore an executor and auditor, not
the primary definition.

Three statements are exact: covariance under the zeta involution, the formula for
the native closure defect, and the continuous theta--Mellin detector identity.
One statement remains open: every non-trivial zero must be shown to occupy the
canonical self-dual PhaseNav shell.  This distinction is maintained throughout.

\section{Symmetric theta channels}

Let
\begin{equation}
 \psi(x)=\sum_{n=1}^{\infty}e^{-\pi n^2x},\qquad x\ge1.
 \label{eq:theta-tail-native}
\end{equation}
The completed zeta function has the symmetric representation
\begin{equation}
 \xi(s)=\frac12+\frac{s(s-1)}2\int_1^\infty
 \psi(x)\left(x^{s/2}+x^{(1-s)/2}\right)\frac{dx}{x}.
 \label{eq:xi-theta-native-x}
\end{equation}
Put $x=e^u$ and write
\begin{equation}
 s=\frac12+z,\qquad z=\delta+it,
 \qquad \delta=\operatorname{Re}s-\frac12.
 \label{eq:centered-z-native}
\end{equation}
Then \cref{eq:xi-theta-native-x} becomes
\begin{equation}
 \xi(s)=\frac12+\frac{s(s-1)}2\int_0^\infty
 \psi(e^u)e^{u/4}
 \left(e^{zu/2}+e^{-zu/2}\right)du.
 \label{eq:xi-theta-native-u}
\end{equation}
The two exponentials are canonical complementary channels; no arbitrary split of
$\xi$ has been introduced.

\section{Eighteen pairs in thirty-six dimensions}

Choose positive quadrature nodes and weights $(u_k,q_k)$, $k=0,\ldots,17$, and
set
\begin{equation}
 b_k=q_k\psi(e^{u_k})e^{u_k/4}>0.
 \label{eq:base-gain-native}
\end{equation}
The native PhaseNav state is the interleaved vector
\begin{equation}
 P_{18}(s)=
 \left(R_{0,+},R_{0,-},\ldots,R_{17,+},R_{17,-}\right)\in\C^{36},
 \label{eq:pn-state}
\end{equation}
where
\begin{equation}
 R_{k,+}(s)=b_ke^{+zu_k/2},\qquad
 R_{k,-}(s)=b_ke^{-zu_k/2}.
 \label{eq:pn-rotor-pair}
\end{equation}
Writing $R_{k,\pm}=\rho_{k,\pm}e^{i\phi_{k,\pm}}$ gives
\begin{align}
 \rho_{k,+}&=b_ke^{+\delta u_k/2}, &
 \phi_{k,+}&=+tu_k/2\pmod{2\pi},\label{eq:pn-plus}\\
 \rho_{k,-}&=b_ke^{-\delta u_k/2}, &
 \phi_{k,-}&=-tu_k/2\pmod{2\pi}.\label{eq:pn-minus}
\end{align}
Thus $t$ acts as phase transport, whereas displacement away from the critical
axis acts as antisymmetric radial gain shear.

The PhaseNav phase grammar is used directly.  The theta seed is an
\texttt{IDENTITY} operator; complex transport is \texttt{DIFFERENTIAL}; the
zeta involution is an \texttt{ORBIT}; and native closure is an
\texttt{AFFECT} operator.

\section{Involution covariance}

Let $X$ swap the two members of every rotor pair and let complex conjugation act
componentwise.

\begin{theorem}[Native involution covariance]
\label{thm:pn-covariance}
For every $s\in\C$,
\begin{equation}
 P_{18}(\J s)=X\overline{P_{18}(s)},
 \qquad \J s=1-\overline{s}.
 \label{eq:pn-covariance}
\end{equation}
\end{theorem}

\begin{proof}
Under $\J$, the centred coordinate satisfies $z\mapsto-\overline z$.  Therefore
\[
 R_{k,+}(\J s)=b_ke^{-\overline z u_k/2}
 =\overline{R_{k,-}(s)},
\]
and similarly $R_{k,-}(\J s)=\overline{R_{k,+}(s)}$.  Applying this identity to
all eighteen pairs proves \cref{eq:pn-covariance}.
\end{proof}

The executable residual is the maximum componentwise norm of the difference
between the two sides.  It is at roundoff level in the regression suite.

\section{The native closure defect}

The canonical self-dual shell is the weighted torus on which every pair has the
base gain $b_k$ in both channels.  Define
\begin{equation}
 \mathcal C_{18}(s)=
 \frac{\sum_{k=0}^{17}
 \left[\log\left(\rho_{k,+}/\rho_{k,-}\right)\right]^2}
 {\sum_{k=0}^{17}u_k^2}.
 \label{eq:pn-closure-defect}
\end{equation}
This quantity is gauge-independent because it uses gain ratios and is unaffected
by any common phase rotation.

\begin{theorem}[Exact native closure identity]
\label{thm:pn-closure}
For all $s\in\C$,
\begin{equation}
 \boxed{\mathcal C_{18}(s)=
 \left(\operatorname{Re}s-\frac12\right)^2.}
 \label{eq:pn-closure-identity}
\end{equation}
Consequently,
\begin{equation}
 \mathcal C_{18}(s)=0
 \quad\Longleftrightarrow\quad
 \operatorname{Re}s=\frac12.
 \label{eq:pn-half-equivalence}
\end{equation}
\end{theorem}

\begin{proof}
From \cref{eq:pn-plus,eq:pn-minus},
\[
 \log\frac{\rho_{k,+}}{\rho_{k,-}}=\delta u_k.
\]
Substitution into \cref{eq:pn-closure-defect} yields
\[
 \mathcal C_{18}(s)=
 \frac{\delta^2\sum_ku_k^2}{\sum_ku_k^2}=\delta^2.
\]
Using \cref{eq:centered-z-native} proves both claims.
\end{proof}

The theorem is independent of known zero ordinates, numerical fitting and the
accuracy of the quadrature.  The same identity holds for every non-degenerate
finite node set and for the continuous state.

\section{The theta--Mellin detector}

The finite detector is
\begin{equation}
 D_{18}(s)=\frac12+\frac{s(s-1)}2
 \sum_{k=0}^{17}\left(R_{k,+}(s)+R_{k,-}(s)\right).
 \label{eq:pn-detector18}
\end{equation}

\begin{proposition}[Continuous detector identity]
\label{prop:pn-detector}
When the quadrature is replaced by the continuous measure in
\cref{eq:xi-theta-native-u}, the detector equals $\xi(s)$ identically.
\end{proposition}

\begin{proof}
Insert the continuous versions of \cref{eq:base-gain-native,eq:pn-rotor-pair}
into \cref{eq:pn-detector18}.  The resulting integral is exactly
\cref{eq:xi-theta-native-u}.
\end{proof}

The finite detector retains the functional equation exactly at the algebraic
level because the two channels are exchanged by $s\mapsto1-s$.  Its numerical
agreement with $\xi$ depends on the finite quadrature and is therefore reported
separately from the exact construction.

\section{The open native-closure bridge}

\begin{postulate}[SOH-PN-C001 / SOH-C004]
\label{post:pn-open-bridge}
Every non-trivial zero $\rho$ of $\xi$ is represented by a zero state that closes
in the canonical self-dual PhaseNav shell:
\begin{equation}
 \xi(\rho)=0\quad\Longrightarrow\quad\mathcal C(\rho)=0.
 \label{eq:pn-open-bridge}
\end{equation}
\end{postulate}

\begin{theorem}[Conditional PhaseNav critical-axis theorem]
\label{thm:pn-conditional-rh}
If Postulate~\ref{post:pn-open-bridge} holds for every non-trivial zero, then every such
zero lies on $\operatorname{Re}s=1/2$.
\end{theorem}

\begin{proof}
Let $\rho$ be a non-trivial zero.  The postulate gives
$\mathcal C(\rho)=0$.  By \cref{thm:pn-closure}, this is equivalent to
$\operatorname{Re}\rho=1/2$.
\end{proof}

This theorem does not hide the hard step.  The functional equation maps an
off-axis zero to a complementary zero, but it does not by itself identify the
two radially distinct states as one native self-dual state.  Establishing
Postulate~\ref{post:pn-open-bridge} requires an additional analytic mechanism.

\section{Possible promotion mechanisms}

The construction turns the vague request for a canonical bridge into a finite
list of mathematically testable promotion routes:
\begin{enumerate}
 \item construct a positive theta-kernel Gram operator whose null vectors are
 exactly the detector-zero states and whose positivity forces equal pair gains;
 \item construct a self-adjoint generator for which the native closed states are
 precisely the zero modes;
 \item embed the paired state into a de Branges space and prove the required
 Hermite--Biehler inequality;
 \item derive the closure condition from Weil positivity;
 \item prove uniqueness of the canonical zero-state representation under the
 involution, excluding two radially distinct partners.
\end{enumerate}
Any one of these would have to be proved without assuming an equivalent form of
RH in disguise.

\section{Executable profile and numerical receipt}

The native program contains eighteen Gauss--Legendre nodes on $u\in[0,2]$.  The
choice gives a low-height computational profile, not a global theorem.  The
executor verifies:
\begin{itemize}
 \item $36$ dimensions and $18$ complementary pairs;
 \item involution covariance;
 \item the exact closure identity at several off-axis points;
 \item the functional equation for the finite detector;
 \item agreement with $\xi(1/2)$ and $\xi(2)=\pi/6$ within the declared finite
 quadrature tolerance;
 \item a small detector residual at the first known zero ordinate.
\end{itemize}
The generated CSV receipt records detector magnitude, closure defect and
covariance residual.  Numerical proximity to known zeros is evidence only for
the implementation of the detector; it is not evidence for
\cref{post:pn-open-bridge}.

\section{Interpretation}

The construction identifies the half-axis as the unique locus on which the
canonical theta pair is both self-dual and free of radial shear.  This is a
stronger statement than observing that the functional equation is symmetric
about one half: it supplies an explicit state, an involution, a detector and a
metric obstruction.

What remains is correspondingly precise.  The research problem is no longer to
explain why one half looks natural.  It is to prove that detector-zero states
cannot inhabit two distinct involution-related radial shells.  PhaseNav has
converted that question into the promotion of one explicit postulate.

\status{The covariance theorem, closure theorem and continuous detector identity
are exact.  The zero-to-native-closure implication is open.  No proof of the
Riemann Hypothesis is claimed.}
